% s-085-p085-m004; original PDF 85, 86
\Aphorism{s-085-p085-m004}{83}
Der Philister färbt ab; alle Werthe, die einmal in der Nähe des Philisters gelegen haben, tragen davon noch die Spuren. Kein Mensch versteht heute unter Religion etwas Anderes als den Glauben an Philistergötter, unter Wissenschaft etwas Anderes als philiströsen Materialismus, unter Lebensführung etwas anderes als die Einengung in Philisterberufe; ja selbst das Genie erkennt es als seinen göttlichen Beruf, den Philister zu ärgern (\Foreign{épater le bourgeois}). Wer heute Ideale wie die eben genannten bekämpft, muss immer erst vorausschicken, ob er das Ideal selbst oder seine philisterhaft angestrichene Aussenseite im Auge hat, ob er den Philosophen oder den Zeitungleser anredet, ob er als freier Geist oder als Liberaler gehört sein will.
\par

